
\documentclass{article}
\usepackage{colortbl}
\usepackage{makecell}
\usepackage{multirow}
\usepackage{supertabular}

\begin{document}

\newcounter{utterance}

\centering \large Interaction Transcript for game `taboo', experiment `high\_en', episode 4 with reka{-}flash{-}3.1{-}t0.0.
\vspace{24pt}

{ \footnotesize  \setcounter{utterance}{1}
\setlength{\tabcolsep}{0pt}
\begin{supertabular}{c@{$\;$}|p{.15\linewidth}@{}p{.15\linewidth}p{.15\linewidth}p{.15\linewidth}p{.15\linewidth}p{.15\linewidth}}
    \# & $\;$A & \multicolumn{4}{c}{Game Master} & $\;\:$B\\
    \hline

    \theutterance \stepcounter{utterance}  
    & & \multicolumn{4}{p{0.6\linewidth}}{
        \cellcolor[rgb]{0.9,0.9,0.9}{
            \makecell[{{p{\linewidth}}}]{
                \texttt{\tiny{[P1$\langle$GM]}}
                \texttt{You are playing a collaborative word guessing game in which you have to describe a target word for another player to guess.} \\
\\ 
\texttt{Rules:} \\
\texttt{(a) You have to reply in the form: CLUE: <some text>. Guesses from the other player will start with GUESS.} \\
\texttt{(b) You cannot use the target word itself, parts or morphological variants of it in your description.} \\
\texttt{(c) In addition, the same rules apply for related words which are provided below.} \\
\\ 
\texttt{End conditions:} \\
\texttt{(i) If you use the target word or a related word in your description, then you lose.} \\
\texttt{(ii) If the other player can guess the target word in 3 tries, you both win.} \\
\\ 
\texttt{Let us start.} \\
\\ 
\texttt{This is the target word that you need to describe and that the other player needs to guess:} \\
\\ 
\texttt{western} \\
\\ 
\texttt{Related words are:} \\
\\ 
\texttt{{-} eastern} \\
\texttt{{-} cowboy} \\
\texttt{{-} movie} \\
\\ 
\texttt{Important: You are under time pressure, give short descriptions that are to the point!} \\
            }
        }
    }
    & \\ \\

    \theutterance \stepcounter{utterance}  
    & \multicolumn{4}{p{0.6\linewidth}}{
        \cellcolor[rgb]{0.9,0.9,0.9}{
            \makecell[{{p{\linewidth}}}]{
                \texttt{\tiny{[P1$\rangle$GM]}}
                \texttt{<reasoning>} \\
\texttt{Let's tackle this. The target word is "western," and I need to come up with a clue that doesn't use any of the related words: eastern, cowboy, movie. Also, I can't use parts or morphological variants of "western."} \\
\\ 
\texttt{First, think about what "western" typically refers to. It's a genre of movies, literature, or music set in the American Old West, involving cowboys, outlaws, etc. But since I can't use "cowboy" or "movie," I need to avoid those terms. Maybe focus on other aspects. The Old West is a key part, but "Old West" itself might be too close. Maybe use synonyms or related concepts. Think about settings: deserts, ranches, saloons, frontier towns. But how to describe that without using forbidden words.} \\
\\ 
\texttt{Alternatively, think of the time period. The late 19th century in the United States. But "United States" might be okay? Wait, the rules say not to use the target word or parts/morphological variants. "Western" is the target, so "west" is part of it. So "western" is the target, so "west" is a morphological part. So I can't use "west" either. Hmm, tricky.} \\
\\ 
\texttt{So need to avoid "western," "eastern," "cowboy," "movie." Also, can't use parts of "western," like "west." So maybe focus on other elements. Westerns often involve horses, spurs, hats, maybe guns. But "cowboy" is related, so maybe avoid terms directly associated with that. Alternatively, think of the genre's characteristics. Westerns are often about conflict, lawlessness, survival. Maybe use a metaphor. For example, "A genre where hats meet horses and outlaws ride." But "hats" and "horses" might be okay. Wait, but "outlaws" is allowed? The related words are eastern, cowboy, movie. "Outlaws" isn't listed, so maybe that's okay. But need to check if "outlaw" is related. Probably not directly, but maybe. Alternatively, think of the setting: vast plains, sunsets, frontier life. How to describe that without forbidden terms.} \\
\\ 
\texttt{Another angle: the term "western" can also refer to the direction, but since "west" is part of it, that's off{-}limits} \\
            }
        }
    }
    & & \\ \\

    \theutterance \stepcounter{utterance}  
    & & & \multicolumn{2}{p{0.3\linewidth}}{
        \cellcolor[rgb]{0.9,0.9,0.9}{
            \makecell[{{p{\linewidth}}}]{
                \texttt{\tiny{[GM$|$GM]}}
                \texttt{abort game} \\
            }
        }
    }
    & & \\ \\

\end{supertabular}
}

\end{document}
